\section*{Resumen ejecutivo}
\addcontentsline{toc}{section}{Resumen ejecutivo}
\markboth{\MakeUppercase{Resumen ejecutivo}}{}

\kicker{Subdocumento 13 · Innovaciones}

Terabyte presenta la cartera obligatoria de cinco innovaciones exigida por el
Capítulo 5 de las Bases Administrativas, una por cada tipo del Artículo
28\textdegree{}.

La cartera no se compuso eligiendo tecnologías. Se compuso \textbf{descartando}: se
contrastó cada candidata contra los 139 requerimientos funcionales y las 21
decisiones del alcance ya definido, y el resultado fue que la mayor parte de
las respuestas evidentes del sector —sistema de citas, reconocimiento óptico
del contenedor, asignación algorítmica de posición, alarmas de frío
escaladas, mensajería estándar, portal de autoservicio, cálculo de emisiones—
\textbf{ya son alcance obligatorio de esta propuesta}. Presentarlas como
innovación habría activado la causal de descarte del Artículo 30\textdegree{}.

Lo que quedó después de ese filtro, y de las restricciones no negociables del
Capítulo 10 del Caso, no son tecnologías sino \textbf{cinco huecos} entre lo
que el terminal necesita como resultado y lo que los requerimientos cubren
como conducta:

\begin{enumerate}[leftmargin=*, itemsep=2pt]
  \item El sistema avisa hacia adentro y nunca hacia afuera: cuando una toma
    refrigerada falla, el dueño de la carga no se entera.
  \item La cita de camión existe, pero el proceso sigue suponiendo que el
    transportista controla cuándo estará lista la carga.
  \item No existe dónde ensayar un cambio: cuatro meses y medio de
    congelamiento, 620 recaladas al año y un patio al 90\,\% en temporada.
  \item El CLIENTE compra el equipamiento de terreno, no tiene quién lo
    mantenga y su proveedor no comparte la consecuencia del indicador
    comprometido con el concedente.
  \item Los requerimientos obligan a \emph{medir} las emisiones. Ninguno
    obliga a \emph{reducirlas}.
\end{enumerate}

Cada innovación se desarrolla en su idea, la tecnología que la sustenta, su
alcance, su forma de implementación y su resultado esperado, conforme al
Formulario T-22, y declara expresamente aquello que requiere investigación
adicional. Al menos una de ellas es verificable durante la marcha blanca de la
Etapa 1, antes del mes 16, conforme al \texttt{RT-26.08}.

\begin{bloquePropuesta}
Ninguna de las cinco es transferible a otro caso del llamado: cada una está
anclada a un hecho medido de este terminal. Y ninguna es indispensable para
cumplir el alcance comprometido: todas degradan a la funcionalidad obligatoria
si no rinden lo esperado.
\end{bloquePropuesta}

\noindent\small\textcolor{terabyteGris}{%
Conforme al Artículo 53\textdegree{} de las Bases Administrativas, este
subdocumento no contiene información económica. La valorización de las cinco
innovaciones en el flujo de caja corresponde al Informe 3.}
